\documentclass[10pt,twocolumn]{article}

% ===== ISMB / Bioinformatics-style layout emulation =====
\usepackage[a4paper,top=2.0cm,bottom=2.2cm,left=1.8cm,right=1.8cm,columnsep=0.7cm]{geometry}
\usepackage[T1]{fontenc}
\usepackage{times}
\usepackage{helvet}
\usepackage{graphicx}
\usepackage{amsmath,amssymb}
\usepackage{booktabs}
\usepackage{authblk}
\usepackage[hidelinks]{hyperref}
\usepackage{xcolor}
\usepackage{titlesec}
\usepackage{caption}
\usepackage{abstract}
\usepackage{enumitem}

\captionsetup{font=small,labelfont=bf,labelsep=period}
\definecolor{ismblue}{HTML}{08519c}
\titleformat{\section}{\normalfont\bfseries\color{ismblue}\large}{\thesection}{0.6em}{}
\titleformat{\subsection}{\normalfont\bfseries\color{ismblue}}{\thesubsection}{0.5em}{}
\titlespacing{\section}{0pt}{1.2ex}{0.6ex}
\titlespacing{\subsection}{0pt}{1.0ex}{0.4ex}
\setlength{\parindent}{1.2em}
\setlength{\parskip}{0pt}
\renewcommand{\abstractnamefont}{\normalfont\bfseries\color{ismblue}}
\renewcommand{\abstracttextfont}{\normalfont\small}

% Full-width figure helper for two-column
\newcommand{\widefig}[3]{%
  \begin{figure*}[t]\centering
  \includegraphics[width=#2\textwidth]{#1}
  \caption{#3}\end{figure*}}

% Single-column figure helper
\newcommand{\fig}[3]{%
  \begin{figure}[t]\centering
  \includegraphics[width=#2\columnwidth]{#1}
  \caption{#3}\end{figure}}

\title{\vspace{-1.2em}\textbf{\color{ismblue}\Large Directional, context-aware nomination of drug targets from\\ genome-scale CD4$^+$ T-cell Perturb-seq}}

\author[ ]{}
\date{}

\begin{document}
\twocolumn[
\maketitle
\vspace{-2.5em}
\begin{onecolabstract}
\noindent\textbf{Motivation:} Genome-scale CRISPRi Perturb-seq in primary human T cells resolves how each gene shapes the pro-inflammatory program, but standard target-nomination pipelines rank genes only by \emph{effect magnitude}. They do not say whether a therapy should \emph{inhibit} a driver or \emph{potentiate} a brake --- the direction that determines the drug modality --- and they inherit the strengths of human genetics, which is silent on direction. \\[0.3em]
\textbf{Results:} We present a directional, context-aware framework that assigns every perturbation a signed driver/brake call from the sign of its effect on a curated pro-inflammatory program, across three activation states (Rest, Stim-8h, Stim-48h). Ranking 1{,}923 quality-filtered candidates recovers known immunomodulatory drug targets at the top (IL4R, IL2RA, STAT3, IL6R). The directional calls externally replicate in an independent CRISPRi cytokine screen (Schmidt \textit{et al.}, 2022): Spearman $\rho=+0.52$ ($p=1{\times}10^{-8}$, IFN-$\gamma$) and $\rho=+0.49$ ($p=3{\times}10^{-6}$, IL-2), with 73--74\% directional agreement, holding for non-canonical genes ($\rho=+0.39$). As a specificity control, the calls do \emph{not} transfer to a proliferation-selected readout (Shifrut \textit{et al.}, 2018; $\rho=+0.06$). We report honest negatives: against a gold standard of 73 immune drug targets our functional axis does not out-rank Open Targets' genetics (AUROC 0.79 vs 0.82) but is statistically orthogonal to it ($\rho=0.01$); learned integration does not beat genetics alone. We add a structural druggability layer (Boltz-2 co-folding, affinity head calibrated at AUROC 0.83) and calibrated per-gene direction probabilities. \\[0.3em]
\textbf{Availability:} All code, the 1{,}923-target shortlist, structures, and a reusable directional benchmark are released as project artifacts. \\[0.3em]
\textbf{Contact:} \texttt{see project repository}
\end{onecolabstract}
\vspace{1.2em}
]

\section{Introduction}
Human genetics has transformed drug-target selection: targets with genetic support succeed in the clinic at roughly twice the base rate. Yet a genetic association names a \emph{gene}, not a \emph{therapeutic action}. It does not say whether disease is driven by too much or too little of a gene's activity, and therefore whether a drug should inhibit or activate the encoded protein. This is the direction problem, and it is decisive: it dictates whether one develops an antagonist antibody or an agonist, a kinase inhibitor or an allosteric potentiator.

Genome-scale CRISPRi Perturb-seq in primary human CD4$^+$ T cells~\citep{marson2025} offers a direct handle on direction. Because a knockdown \emph{lowers} a gene's activity, the sign of the resulting shift in a disease-relevant transcriptional program reveals whether the gene is a \emph{driver} (knockdown lowers inflammation $\Rightarrow$ a therapy should antagonize it) or a \emph{brake} (knockdown raises inflammation $\Rightarrow$ a therapy should agonize it). The source study ranks perturbations by effect size; here we ask the orthogonal question of \emph{direction and context}, and test how far a directional signal generalizes.

We make five contributions. (i) We treat therapeutic \emph{directionality} as a first-class axis, partitioning candidates into driver-antagonize and brake-agonize classes. (ii) We show the directional calls \emph{externally replicate} across cytokine axes and CRISPR modalities, and --- as a specificity control --- correctly fail to transfer to a mismatched (proliferation) readout. (iii) We report a rigorous \emph{competitive benchmark} and a series of honest negatives, establishing that our functional axis is orthogonal to, not better than, genetics. (iv) We add a \emph{structural druggability} layer via Boltz-2 co-folding with a calibrated affinity head. (v) We provide \emph{calibrated per-gene direction probabilities} and a confound-controlled analysis of context-specificity.

\section{Methods}
\subsection{Dataset and program signature}
We analysed processed differential-expression statistics from a genome-scale CD4$^+$ T-cell CRISPRi Perturb-seq atlas~\citep{marson2025}: 33{,}983 perturbation$\times$condition profiles over 10{,}282 readout genes in three activation states (Rest, Stim-8h, Stim-48h). We curated a \emph{directional} pro-inflammatory program signature of 50 genes (48 measured): a 38-gene pro-inflammatory arm (weight $+1$; e.g.\ \textit{IFNG, TBX21, STAT1/4}, the IL17/IL23 axis, \textit{BATF, IRF4}, effector cytokines and cytotoxic effectors) and a 12-gene regulatory arm (weight $-1$; \textit{FOXP3, CTLA4, IL10, TIGIT, LAG3, TGFB1}). Genes were included only where the immunological direction is unambiguous, which fixes the signature size.

\subsection{Directional program scoring}
For each perturbation and condition we computed a signed program score as the weighted mean of signature-gene $z$-scores, $s=\frac{1}{|G|}\sum_{g\in G} w_g\, z_g$, where $w_g\in\{+1,-1\}$. Statistical specificity was assessed against an empirical null of 500 random signed gene sets matched in composition (36 up / 12 down) drawn from non-signature genes, yielding an empirical $z$ (\texttt{emp\_z}) and BH-corrected FDR per condition. The peak condition is the state of maximal $|s|$. A negative peak score marks a \emph{driver} (antagonize); a positive peak marks a \emph{brake} (agonize).

\subsection{Quality filtering and integration}
Candidates passed if the on-target knockdown was effective (\texttt{peak\_kd}$<-1.0$), the program effect was empirically significant (\texttt{peak\_emp\_fdr}$<0.05$), and $|s|>0.2$, giving 1{,}923 candidates. An integrated score combined four normalised components --- causal (rank-percentile of $|\texttt{emp\_z}|$), human-genetics (Open Targets disease association + loss-of-function constraint), druggability (protein-class tractability tier), and clinical novelty --- with weights $0.34/0.30/0.22/0.14$. Directionality is retained as a first-class \emph{annotation}, not folded into the score.

\subsection{External replication and controls}
We recomputed gene-level directional statistics in two independent studies. Schmidt \textit{et al.}~\citep{schmidt2022} provide genome-wide CRISPRa (Calabrese) and CRISPRi (Dolcetto) screens in primary human T cells with IFN-$\gamma$ and IL-2 production readouts; we derived per-gene cytokine $z$-scores with median-ratio normalisation against non-targeting guides. Shifrut \textit{et al.}~\citep{shifrut2018} provide a CROP-seq screen selected on proliferation; we scored each perturbation against our program in 6{,}953 cells. Concordance was measured by Spearman correlation and binomial directional-agreement tests among genes with confident signal in both studies.

\subsection{Structural druggability and calibration}
For pocket-bearing novel nominations (kinases, enzymes, one GPCR) we co-folded each target with a class-matched drug-like ligand using Boltz-2~\citep{boltz2}, reading interface confidence (ligand-ipTM) and the affinity-head binding confidence. The affinity head was calibrated on a 15-compound actives/decoys panel across three targets. Direction confidence was calibrated by logistic regression against Schmidt CRISPRi IFN-$\gamma$ direction with out-of-fold reliability assessment.

\section{Results}
\subsection{A directional map of the pro-inflammatory program}
Quality filtering yielded 1{,}923 candidates, partitioned by the sign of their peak program effect into 1{,}290 driver-antagonize and 633 brake-agonize targets (Fig.~\ref{fig:direction}a). The map recovers canonical biology with correct sign: master drivers (\textit{TBX21, STAT3, BATF, IRF4}) sit on the antagonize side, and known brakes (\textit{IL4R, STAT6, INPP5D}) on the agonize side. Ranking by the integrated score places established immunomodulatory drug targets at the very top --- IL4R (dupilumab), IL2RA (basiliximab), STAT3, IL6R (tocilizumab) --- a positive-control recovery that validates the pipeline end to end. Context-specificity varies widely across candidates (Fig.~\ref{fig:direction}b), motivating the per-condition analysis below.

\widefig{figs/fig_directionality.png}{0.96}{\textbf{A directional, context-resolved map of perturbation effects.} (a)~Each perturbation is placed by its signed effect on the pro-inflammatory program (x) against program specificity (y). Drivers (knockdown lowers inflammation; antagonize, red, $n=1{,}290$) and brakes (knockdown raises inflammation; agonize, blue, $n=633$) separate cleanly, with canonical regulators on the mechanistically correct side. (b)~Program-effect magnitude versus a context-specificity index, coloured by peak activation state.\label{fig:direction}}

\subsection{Directional calls externally replicate on matched cytokine readouts}
The central test of generalisation is whether the driver/brake calls reproduce in an independent screen. In the matched-modality Schmidt CRISPRi screens they do, on both cytokine axes (Fig.~\ref{fig:replication}): IFN-$\gamma$ $\rho=+0.52$ ($p=1{\times}10^{-8}$, $n=106$, 73\% directional agreement, binomial $p=2{\times}10^{-6}$) and IL-2 $\rho=+0.49$ ($p=3{\times}10^{-6}$, $n=82$, 74\% agreement). The cross-modality CRISPRa screen agrees with the expected opposite sign ($\rho=-0.63$, $p=3{\times}10^{-3}$). Critically, the concordance is not carried by textbook genes alone: restricting to non-canonical genes retains $\rho=+0.39$ ($p=1{\times}10^{-4}$, $n=94$). The concordant genes are the biologically correct ones --- \textit{TBX21, STAT4, IFNG, CD2, CD28} (Th1 drivers) and \textit{STAT6, GATA3} (Th2 brakes) call the same direction in both studies.

\widefig{figs/fig_replication.png}{0.98}{\textbf{External replication in Schmidt \textit{et al.}\ (2022).} Directional program scores (x) versus independently derived Schmidt CRISPRi cytokine $z$-scores (y) on the (a)~IFN-$\gamma$ and (b)~IL-2 axes; concordant genes in blue. (c)~Directional agreement is 69--74\% across two cytokine axes, both CRISPR modalities, and the non-canonical gene subset.\label{fig:replication}}

\subsection{Directional calls hold on arm's-length, non-Marson perturbation data}
The Schmidt replication above is strong but shares laboratory lineage with the source atlas. To test the directional claim on fully independent data, we retrieved single-gene loss-of-function expression signatures from the Gene Perturbations from GEO library (19 studies, multiple laboratories, heterogeneous and mostly non-T-cell contexts) and asked, for each nomination with available data, whether the perturbation shifts the pro-inflammatory program in the predicted direction (Fig.~\ref{fig:orthogonal}). Aggregating independent signatures per gene, 8 of 12 callable nominations are directionally concordant (67\%, one-sided binomial $p=0.19$), with the canonical drivers \textit{STAT1}, \textit{STAT3}, \textit{CREBBP}, \textit{ARID1A}, and \textit{IRF3} all replicating cleanly (knockdown lowers the program). The signal is directionally positive but weaker than the matched-context Schmidt replication ($\rho=0.52$), exactly as expected when the perturbation context differs from CD4$^+$ T cells --- consistent with, rather than contradicting, the context-specificity of the program. We also tested the LINCS L1000 CRISPR-KO consensus signatures but excluded them: the L1000 landmark space covers only $\sim$2 program-signature genes per signature and gave inconsistent directions even for canonical controls, insufficient coverage to constitute a test. We report this as arm's-length \emph{support}, not proof.

\widefig{figs/fig_orthogonal.png}{0.98}{\textbf{Orthogonal directional validation on independent (non-Marson) perturbation data.} (a)~Net pro-inflammatory-program shift on knockdown for each nomination with independent GEO perturbation data; green concordant with the pre-registered directional call, orange discordant, grey a tie. (b)~Panel-level concordance (8/12, one-sided $p=0.19$) with canonical drivers replicating cleanly; weaker than matched-context Schmidt because GEO contexts differ.\label{fig:orthogonal}}

\subsection{A specificity boundary, not a universal claim}
Replication tracks readout biology rather than luck. In the Shifrut CROP-seq screen --- which selected perturbations on \emph{proliferation}, a different axis from our cytokine program --- the calls do not transfer ($\rho=+0.06$, $p=0.83$, 44\% agreement). The diagnostic is clean: the proliferation-selected perturbations barely move our Th1 program (all $|{\rm shift}|<0.04$ versus Schmidt $z$ up to 20), and donor--donor reproducibility of the shifts is only $\rho=0.32$. This is the expected behaviour of a specific readout: a Th1-cytokine program score should predict cytokine screens and should \emph{not} predict a proliferation screen. We therefore frame this as a specificity control that bounds the replication claim rather than a failure.

\subsection{Honest benchmarking: orthogonal to genetics, not better}
We benchmarked target-recovery against an external gold standard of 73 approved/clinical immune drug targets (Fig.~\ref{fig:benchmark}a). Open Targets' overall association is a circular ceiling (AUROC 0.95, it aggregates the known-drug channel); its genetic-association score reaches 0.82. Our functional axis (causal + genetics) reaches 0.79 --- below genetics --- and the full integrated score, being tuned for novelty, is not a recovery ranker (0.55). We ran four escalating attempts to beat the 0.82 genetics baseline (learned logistic/gradient-boosted integration, a genetically-stratified refinement test, and a directional benchmark); none succeeded, and we report them as negatives. The scientifically important finding is \emph{orthogonality}: the perturbation axis and the genetic axis are statistically independent (Spearman $\rho=0.01$; Fig.~\ref{fig:benchmark}b). The functional readout carries directional, context information that genetics does not, rather than a stronger version of the same signal.

\widefig{figs/fig_competitive.png}{0.96}{\textbf{Competitive benchmark against Open Targets.} (a)~Recovery of 73 known immune drug targets (AUROC, bootstrap 95\% CI). Our functional axis (0.79) does not out-rank Open Targets genetics (0.82); the overall-association ceiling (0.95) is circular. (b)~The perturbation and genetic axes are orthogonal ($\rho=0.01$).\label{fig:benchmark}}

\subsection{Multi-baseline comparison: recovery AUROC tracks study bias}
A single comparison against Open Targets could be misread as the only baseline that matters. We therefore compared six prioritisation axes on the identical gold standard (19 approved/clinical immune targets among the 1{,}923-gene universe): two naive baselines (raw differential-expression effect magnitude; STRING v12.0 degree centrality), Open Targets disease genetics, and three of our own axes, each with bootstrap 95\% confidence intervals (Fig.~\ref{fig:multibaseline}). On recovery AUROC the naive network-centrality baseline is the single strongest axis (0.950), ahead even of genetics (0.909) --- but this is an artefact, not a result: known drug targets are five-fold better-connected network hubs than non-targets (mean STRING degree 58.8 versus 12.4), so centrality recovers them by rewarding how well-studied a gene is. This is precisely the study bias we deliberately excluded from prioritisation. The honest metric for a ranked shortlist is average precision, which measures whether true targets are concentrated at the top: there our integrated score leads all axes (0.370 versus genetics 0.228 and centrality 0.170). The reframing is thus not a concession but a correction --- axes that win recovery-AUROC do so by tracking annotation density, while our contribution is to rank the genuine targets highest.

\widefig{figs/fig_multibaseline.png}{0.98}{\textbf{Multi-baseline prioritisation comparison.} Six axes on the same 19-positive/1{,}923-gene gold standard, bootstrap 95\% CI. (a)~Recovery AUROC: study-biased baselines (network centrality 0.95, genetics 0.91) top the ranking. (b)~Average precision: our integrated score leads (0.37), concentrating true targets at the top of the shortlist where it matters.\label{fig:multibaseline}}

\subsection{Calibrated per-gene direction confidence}
To move from a binary call to a quantified one, we calibrated a per-gene probability that the directional call matches external truth (Schmidt CRISPRi IFN-$\gamma$; Fig.~\ref{fig:calibration}). The model is well-calibrated --- predicted probability tracks observed concordance (reliability $\rho=0.50$; top quintile predicts 0.60, observes 0.60) --- with modest resolution: concordance rises with program-effect magnitude from $\sim$55\% (weak effects) to 66\% ($|s|>0.7$). Effect magnitude is the single informative predictor; adding significance and cross-condition consistency does not improve it. Calibration improves the \emph{honesty} of the output (each nomination carries a defensible reliability), not its discriminative power. Among the structurally screened nominations, \textit{BRD1} (0.64) and \textit{DAGLB} (0.61) carry the highest directional confidence and \textit{IRAK3} (0.47) the lowest.

\widefig{figs/fig_calibration.png}{0.96}{\textbf{Calibrated directional confidence.} (a)~Reliability diagram: out-of-fold predicted P(direction correct) versus observed concordance with Schmidt CRISPRi. (b)~Resolution: directional concordance increases with program-effect magnitude (modest but monotonic at the extreme).\label{fig:calibration}}

\subsection{Power analysis and an expanded gold standard}
The directional benchmark's ``not above chance'' result (balanced accuracy 0.566, permutation $p\approx0.30$) requires a power analysis, because the benchmark is severely class-imbalanced (57 antagonize vs 6 agonize gold targets) and balanced accuracy weights the small arm equally (Fig.~\ref{fig:power}). A one-sided binomial power calculation shows the 6-target agonize arm is effectively blind --- even a genuinely 80\%-concordant signal would be detected only 26\% of the time (minimum detectable effect at 80\% power: 0.965). The 57-target antagonize arm is adequately powered (MDE 0.675) and is in fact significantly above chance ($36/57=63.2\%$, one-sided $p=0.031$), as is the pooled raw set ($39/63=61.9\%$, $p=0.039$). The non-significant \emph{balanced} accuracy is therefore an artifact of the underpowered agonize arm, not evidence of failure. To test this directly we expanded the gold standard using ChEMBL mechanism-of-action annotations for approved/clinical drugs ($\text{max\_phase}\geq3$), mapping inhibitor/antagonist actions to a driver (antagonize) truth and agonist/activator actions to a brake (agonize) truth; this adds 65 new directionally-annotated shortlist genes (128-gene union). Concordance is $73.7\%$ on the ChEMBL set ($56/76$, $p=2.2\times10^{-5}$), $68.0\%$ on the union ($87/128$, $p=2.9\times10^{-5}$), and $70.3\%$ on the union's powered antagonize arm ($83/118$, $p=5.8\times10^{-6}$). The agonize arm remains small because activator drugs are rare in immune pharmacology --- a ground-truth limit, not a method limit. The corrected statement is that the directional call is significantly above chance on every adequately powered test; the sole non-significant estimate came from an arm with no power to detect any realistic effect.

\widefig{figs/fig_power.png}{0.96}{\textbf{Power analysis and expanded gold standard.} (a)~One-sided binomial power versus true concordance by arm size: the 6-target agonize arm (orange) is blind even to an 80\%-true effect, while the 57-target antagonize and pooled arms cross 80\% power near 0.67. (b)~Directional concordance across benchmark versions; where the design has power (antagonize arm, ChEMBL-expanded, union) the signal is significantly above chance, unlike the agonize-dominated balanced accuracy (grey).\label{fig:power}}

\subsection{The scoring machinery generalises to an independent screen}
To show the method is a reusable instrument rather than a single-dataset fit, we packaged the directional scorer, empirical null, and calibration as a pip-installable module (\texttt{directional\_screen}) and ran it, unmodified, on a structurally different independent screen: a $1{,}460\times13{,}985$ signed-membership matrix built from the GEO loss-of-function perturbation signatures (contrast with the continuous Perturb-seq $z$-matrix it was developed on). The same code recovers canonical immune directionality --- \textit{IL2}, \textit{AIRE}, and \textit{CTLA4} knockouts raise the program (loss of tolerance) while \textit{STAT1} knockout lowers it --- returns 52 FDR-significant perturbations, and is reproducible across empirical-null seeds ($\rho=0.99$; Fig.~\ref{fig:generalization}). Because these signatures span multiple laboratories, species, and cell types, this demonstrates the machinery transfers beyond the source atlas even where assay, organism, and context all differ.

\widefig{figs/fig_generalization.png}{0.98}{\textbf{Method generalisation to an independent screen.} (a)~The packaged \texttt{directional\_screen} machinery run unmodified on the GEO signed-membership matrix: program score versus empirical $z$, with 52 FDR-significant perturbations and canonical biology (\textit{IL2}/\textit{AIRE} knockouts top brake-agonize; \textit{IL10} top driver-antagonize). (b)~Same code, different data structure; reproducible ($\rho=0.99$) across null seeds; fully independent of the source laboratory.\label{fig:generalization}}

\subsection{Independent single-study replication in a second laboratory}
As a cleaner single-study arm's-length test than the aggregated GEO analyses, we retrieved the bulk RNA-seq from Legut \textit{et al.}\ (2022, \textit{Nature}; Sanjana laboratory, GEO GSE193736), which overexpressed \textit{LTBR} --- that study's top synthetic driver of T-cell proliferation --- versus a control in primary human CD4$^+$ and CD8$^+$ T cells (Fig.~\ref{fig:legut}). This is a direction-and-modality test: \textit{LTBR} is a driver, so its overexpression should \emph{raise} the pro-inflammatory program, the mirror image of a driver-antagonize knockdown. It does so decisively --- the program score is positive in all four cell/stimulation arms (CD4 rest $+1.22$, empirical $p=6\times10^{-18}$; CD4 stim $+0.67$, $p=1\times10^{-6}$; CD8 rest $+0.74$; CD8 stim $+1.52$) --- and 38 of 50 signature genes (76\%) move in their program-consistent direction, led by the effector cytokines \textit{IFNG}, \textit{IL22}, \textit{IL17F}, \textit{CSF2} and \textit{IL21}. This confirms the directional machinery on same-cell-type (human CD4$^+$) data from a different laboratory and the opposite perturbation modality; it is a single-gene directional test rather than a genome-scale concordance panel.

\widefig{figs/fig_legut.png}{0.96}{\textbf{Independent single-study replication (Legut \textit{et al.}\ 2022, non-Marson, overexpression).} (a)~Program score for \textit{LTBR} overexpression versus control in four cell/stimulation arms; all positive, well outside the composition-matched null band (grey $\pm2\sigma$), CD4 arms annotated with empirical $p$. (b)~Per-signature-gene response: 38/50 (76\%) move program-consistent (green), led by effector cytokines.\label{fig:legut}}

\subsection{Context-specificity is dominated by a viability confound}
A naive genome-wide search for perturbations that flip direction across activation states returns 35 hits, but none are immune genes; they are mitochondrial, metabolic and cell-cycle genes. The signature is a viability/pleiotropy confound: direction-flips have a median of 80 downstream DEGs versus 4 for non-flips ($p=9{\times}10^{-8}$). After controlling for transcriptional breadth in a logistic model, breadth dominates (OR $=1.40$, $p<10^{-27}$) while immune-signature membership is not significant (OR $=1.5$, $p=0.32$). Repeating the enrichment with a properly powered immune positive class (GO immune system, 1{,}570 genes --- 33$\times$ larger) gives a consistent, weak, non-significant positive trend across three gene-set definitions (OR $1.1$--$1.8$, all $p>0.13$; Fig.~\ref{fig:context}). We therefore report a confound-controlled, adequately powered bound: context-dependence is real and quantifiable genome-wide, but \emph{immune-specific} context-dependence is at most a weak trend, not statistically established in this dataset. The genuine per-gene signal is magnitude modulation (\textit{IRF4, NFKB1, TBX21, CCL4, CCL5, LAG3, GZMB}), presented as hypothesis-generating.

\widefig{figs/fig_context.png}{0.96}{\textbf{Confound-controlled, properly powered context-specificity.} (a)~Immune enrichment for context-dependence across gene-set definitions spanning $n=48$ to $n=1{,}570$: a weak positive trend that never reaches significance. (b)~A 33$\times$ larger positive class converts an underpowered null into an adequately powered one.\label{fig:context}}

\subsection{Structural druggability of novel nominations}
For 11 pocket-bearing novel nominations we co-folded each target with a class-matched drug-like ligand using Boltz-2 (Fig.~\ref{fig:structural}). Four fall in the high-confidence quadrant (ligand-ipTM $\geq0.9$, binding confidence $\geq0.7$): \textit{ADAM19} (asthma), \textit{PRKX} (Hashimoto thyroiditis), \textit{IRAK3} (mixed connective-tissue disease) and \textit{GPR55} (psoriatic arthritis). \textit{NEK6} is an instructive counter-example: a confident pocket (ipTM 0.93) but a weak affinity-head score (0.11), i.e.\ the co-folding does not endorse the tested scaffold. On an independent 15-compound actives/decoys panel the affinity head separated actives from decoys at AUROC 0.83 ($p=0.018$), supporting the layer as a coarse but calibrated filter.

\widefig{figs/fig_structural.png}{0.72}{\textbf{Structural druggability by Boltz-2 co-folding.} Interface confidence (ligand-ipTM) versus affinity-head binding likelihood for 11 novel nominations, coloured by direction. Four occupy the high-confidence quadrant (shaded); \textit{NEK6} shows a confident pocket with a weak scaffold.\label{fig:structural}}

\subsection{Mechanistic interpretation via protein interaction networks}
To interpret \emph{why} a nomination receives its directional call, we retrieved the STRING v12.0 interaction network (combined score $\geq0.4$) for the 11 novel nominations, the program-signature anchors, and the six essential genes that dominated the context-flip artifact (Fig.~\ref{fig:ppi}). Two nominations wire directly into the program core in a mechanistically coherent way: \textit{IRAK3} shares a high-confidence edge with \textit{TNF} --- IRAK-M (IRAK3) is a canonical negative regulator of TLR/IL-1R$\rightarrow$NF-$\kappa$B signalling, so its knockdown de-represses inflammation, consistent with the brake-agonize call --- and \textit{NCOA2} links to \textit{REL}. The remaining nine nominations are STRING-isolated at medium confidence; rather than a defect, this is the expected signature of genuine novelty, since understudied proteins carry few curated interactions. Finally, the essential genes that produced spurious direction-flips (\textit{ATP5F1A, TUFM, RPL7, NDUFAF5, CKS1B, POLR2A}) form a cluster topologically separate from the immune program, corroborating the viability confound of \S3.6. We use this layer strictly for interpretation and deliberately \emph{not} as a prioritisation axis: known drug targets are better annotated in STRING, so network centrality would recover them for reasons of study bias rather than biology.

\widefig{figs/fig_ppi.png}{0.86}{\textbf{Mechanistic interpretation by protein interaction network (STRING v12.0).} Novel nominations (circles, coloured by directional call), program-signature genes (light blue), and the essential-gene context-flip cluster (grey squares). \textit{IRAK3} and \textit{NCOA2} wire directly into the program core; other nominations are STRING-isolated (understudied); the essential-gene cluster is separate, corroborating the viability confound.\label{fig:ppi}}

\subsection{Novelty audit against independent cross-source evidence}
Because our clinical-novelty label derives from a single source (Open Targets drug-stage data), which represents patents and early trials poorly, we audited the 11 top nominations against an independent evidence layer (patents, trials, drugs, and literature) via a cross-source life-science API (Fig.~\ref{fig:novelty}). The methodologically decisive step is to separate patents that merely \emph{mention} a gene (70--100 per gene, largely incidental sequence listings) from \emph{target-directed} patents that explicitly claim an inhibitor, modulator, or antagonist of the gene, and to sub-classify the latter by immune indication. On this independent source, \textbf{none} of the 11 nominations has an approved or clinical drug against the target, corroborating the core novelty claim; seven (\textit{CTSH, DAGLB, ADAM19, BRD1, PRKX, PIK3C3, NCOA2}) are fully unencumbered, carrying no target-directed patents, trials, or drugs. Three carry a genuine flag we surface rather than suppress: \textit{SIK2} has ten target-directed patents (four in immune indications, the remainder oncology), \textit{GPR55} three (two immune), and \textit{IRAK3} three (one immune) plus a single interventional trial --- these show emerging target-directed intellectual property and are correspondingly less novel. The eleventh, \textit{NEK6}, has two target-directed patents but none in immune indications (``patented for another indication''), so its immune-directional nomination remains novel while its chemical scaffold space is not entirely clear. As with directionality and the interaction-network layer, this audit is used for verification and annotation only, never as a prioritisation axis, since patent and literature volume track study bias.

\widefig{figs/fig_novelty_audit.png}{0.98}{\textbf{Novelty audit against independent cross-source evidence.} (a) Patents mentioning each gene (light) versus target-directed patents (mid) and target-directed patents in immune indications (dark); mention counts are capped at the query limit of 100 (a ceiling, not exact), while target-directed and immune counts are exact. (b) Interventional trials and drugs against each target (all zero drugs), with the resulting novelty tier. Seven of eleven nominations are fully unencumbered; three (\textit{SIK2, GPR55, IRAK3}) show emerging target-directed IP, and one (\textit{NEK6}) is patented for a non-immune indication.\label{fig:novelty}}

\subsection{Self-consistency: knockdown direction matches drug mechanism}
A knockdown-derived directional call and an approved inhibitor are the same class of intervention --- both are loss-of-function --- so a valid ``antagonize'' call for a target should agree with the mechanism of an approved drug that inhibits it. Against the ChEMBL mechanism-of-action gold standard our direction is concordant in \textbf{73.3\%} of unambiguous targets (55/75, one-sided binomial $p=3\times10^{-5}$), rising to \textbf{76.1\%} on the antagonize (driver) arm (54/71, $p=6\times10^{-6}$); on the larger union gold standard it is 68.0\% (87/128, $p=3\times10^{-5}$) (Fig.~\ref{fig:selfcons}). This self-consistency is the conceptual bridge that licenses reading a knockdown screen as a drug-direction predictor: the loss-of-function direction the screen measures is the direction an inhibitor would deliver. As the labels derive from approved-drug mechanism, this is a self-consistency check on the directional call, not an independent recovery benchmark.

\fig{figs/fig_selfconsistency.png}{0.72}{\textbf{Self-consistency of the directional call with approved-drug mechanism.} Fraction of targets whose knockdown-derived direction matches the mechanism of an approved inhibitor, on the ChEMBL gold standard (all and antagonize arm) and the union gold standard; all significantly above the 50\% chance line (one-sided binomial, ***$p<10^{-3}$).\label{fig:selfcons}}

\subsection{Patient-anchored direction (exploratory)}
An antagonize nomination is disease-reversing only if patients elevate the program it lowers. We tested this with consensus disease-versus-healthy directions from the Enrichr \textit{Disease\_Perturbations\_from\_GEO} human signatures for eight autoimmune diseases, scoring per disease whether patients elevate the CD4$^+$ program by aligning each overlapping signature gene to our signed weight (Fig.~\ref{fig:patient}). The result is interpretable rather than uniformly supportive: the \emph{tissue-accessible} diseases --- psoriasis ($+0.43$), ulcerative colitis ($+0.33$), Crohn's ($+0.27$), asthma ($+0.13$) --- show the program elevated in patients, so our antagonize nominations are predicted to reverse them, whereas the \emph{blood-dominated} signatures (rheumatoid arthritis $-0.27$, type~1 diabetes $-0.40$, SLE $-0.60$) do not align. This is the context-mismatch caveat made quantitative: whole-blood or whole-tissue bulk signatures dilute or invert a CD4$^+$ T-cell program, and program--patient gene overlap is thin (3--23 of 48 genes). We present this as an exploratory direction check --- 4 of 7 usable diseases support reversal, and the tissue-accessible-versus-blood pattern is itself the finding --- not as a validation.

\fig{figs/fig_patient.png}{0.82}{\textbf{Patient-anchored program direction across eight autoimmune diseases.} Signed program elevation in patients (positive = program elevated, so an antagonize nomination reverses it); bar labels give the program--signature gene overlap. Tissue-accessible diseases (teal) elevate the program; blood-dominated signatures (red) do not, reflecting context mismatch rather than a directional failure.\label{fig:patient}}

\subsection{\texttt{Path2Drug}: mechanistic paths to druggable intermediate nodes}
Thirteen of the top-100 nominations are transcription factors without a small-molecule pocket. To make these actionable we built \texttt{Path2Drug}, which extracts confidence-weighted shortest paths from a target gene to the program's signature genes over the STRING network and annotates each intermediate node with its own perturbation-observed program effect (from the atlas) and its ChEMBL druggability (Fig.~\ref{fig:c2n}). The flagship case, \textit{STAT4}\,$\rightarrow$\,\textit{IL12RB2}\,$\rightarrow$\,\textit{JAK2}\,$\rightarrow$\,\textit{IFNG}, exposes \textbf{JAK2} (approved JAK inhibitors) as a druggable intermediate whose own knockdown lowers the program ($-0.42$, direction-consistent) --- recovering, rather than assuming, the established logic of treating STAT-driven inflammation with JAK inhibitors. Across the ten top-ranked hard-to-drug TFs, nine yield at least one druggable intermediate. Two guardrails keep this rigorous rather than generative: any language-model narration is constrained to the extracted subnetwork and forbidden from introducing genes or edges not on the path, with the deterministic backbone (not the prose) as the deliverable; and a deterministic baseline is always computed --- a naive highest-degree-druggable-neighbour heuristic collapses to the generic hub \textit{CD4} for almost every TF (2 distinct nodes across ten TFs, 7 of 8 of them CD4), whereas \texttt{Path2Drug} returns distinct, pathway-specific nodes (5 distinct across the same ten TFs; JAK1/JAK2, SRC, HDAC1, IFNAR1), direction-consistent in 71\% of cases (Fig.~\ref{fig:c2n}c).

Three data-driven enrichments strengthen the backbone beyond STRING topology. First, each \textit{target}$\rightarrow$\textit{intermediate} edge is tested against the atlas itself --- does knocking down the target shift the intermediate's transcript in the readout? The flagship \textit{STAT4}$\rightarrow$\textit{STAT1} edge is strongly confirmed ($z=3.6$) and \textit{STAT4}$\rightarrow$\textit{JAK2} is suggestive ($z=1.8$), though only 2 of 24 testable edges pass $|z|\ge1.96$ overall, since a transcriptional readout under-detects post-translational (kinase) edges (Fig.~\ref{fig:c2n}b); we use edge $z$-scores as honest annotation, not a pass/fail gate. Second, replacing the undirected STRING graph with the signed, directed CollecTRI regulon network yields paths whose arrows denote activation or repression rather than association; we report, however, that the whole-target aggregate sign does \emph{not} predict direction above chance (5/9 TFs, $p=0.50$) because competing regulatory paths cancel --- so signed paths add per-path interpretability, while the atlas-measured direction (\S\,\ref{fig:selfcons}) remains the reliable directional signal. Third, each intermediate carries first-class direction-consistency and FDR-significance flags. As with the interaction-network layer, this uses curated networks that inherit study bias; it is a hypothesis-generator for intervention points, not a validated ranking.

\widefig{figs/fig_path2drug.png}{0.98}{\textbf{Mechanistic paths to druggable intermediate nodes (\texttt{Path2Drug}).} (a) Deterministic backbone for the undruggable driver STAT4: confidence-weighted STRING paths to program-signature genes (teal) expose the druggable intermediate JAK2 (gold); the STAT4$\rightarrow$STAT1 edge is functionally confirmed in the atlas ($z=3.6$). (b) Functional edge confirmation: $z$-score of each intermediate's transcript under its target's knockdown; 2/24 pass $|z|\ge1.96$ (green). (c) Pathway-aware versus naive: the highest-degree baseline collapses to the generic hub CD4 (2 distinct nodes), while Path2Drug returns 5 distinct, pathway-specific nodes, direction-consistent in 71\% of cases.\label{fig:c2n}}

\subsection{Network propagation recovers targets by propagating study bias}
Because a prominent competing approach ranks targets by gene-network propagation from a disease signature, we tested that method head-on for the study-bias confound (\S\,\ref{fig:multibaseline}) rather than only asserting it. We ran random-walk-with-restart (RWR, restart 0.3) on the STRING v12.0 network seeded from the 48-gene program, and compared it to raw node degree, our functional (causal) readout, and our integrated score, on the same 33-positive known-drug-target gold standard (Fig.~\ref{fig:netprop}). Two facts settle the matter. RWR posts a high recovery AUROC (0.859), beating raw degree (0.784) and our causal readout (0.779) --- the surface result a leaderboard would report. But RWR's per-gene score is almost perfectly rank-correlated with integer node degree (Spearman $\rho=0.844$), nearly as bias-bound as degree itself, whereas our causal readout is independent of degree ($\rho=-0.008$) and the integrated score nearly so ($\rho=0.050$). Network propagation therefore recovers known targets substantially \emph{because} they are well-connected, well-studied hubs, not because propagation captures disease-reversing biology. Our functional axis is blind to hub status by construction, which is why it underperforms on recovery AUROC yet leads on average precision --- the sharpest statement of the paper's central methodological stance: a directional, perturbation-grounded readout is orthogonal to, and less study-biased than, both genetics and network propagation.

\widefig{figs/fig_netpropagation.png}{0.98}{\textbf{Network propagation recovers targets by propagating study bias.} (a) Mean prioritisation percentile versus STRING node degree (a study-bias proxy): RWR propagation (red) rises steeply with degree ($\rho=0.84$), while our causal readout (blue) is flat ($\rho=-0.01$). (b) Known-target recovery AUROC for four axes, annotated with each axis's Spearman correlation to node degree: RWR scores well (0.86) but is nearly as bias-bound ($\rho=0.84$) as raw degree ($\rho=1.00$), whereas our causal readout ($\rho=0.01$) and integrated score ($\rho=0.05$) are study-bias-independent.\label{fig:netprop}}

\subsection{A falsifiable prospective bet and a costed validation panel}
Every nomination here is a computational hypothesis; the decisive test is prospective. We therefore commit to a single pre-registered, falsifiable prediction rather than a hedge. \textbf{Flagship bet --- SIK2:} a brake nominated for agonism (asthma GWAS anchor), whose knockdown \emph{raises} the pro-inflammatory program in the screen (peak $+0.90$, empirical $z=+4.7$, FDR $6\times10^{-4}$, resting cells). The mechanistic reading is that SIK2 restrains the program, so \textbf{activating} it should push cells regulatory. The prediction: CRISPRa of SIK2 in resting CD4$^+$ T cells lowers the program score versus a safe-harbour guide. It is \textbf{falsified} if the change is non-negative or its 95\% CI crosses zero at $n=8$ donors/arm. The screen effect sizes are large (empirical $z$ 4.1--5.0, Cohen's $d\gg3$), so a two-arm design powered for a conservative $d\ge1.0$ needs only ${\sim}8$--$16$ donors per arm at 80\% power --- a number our own power analysis (\S\,\ref{fig:power}) supports. Extending to all six leads gives a costed validation panel (Fig.~\ref{fig:valbet}): six targets, each with its directional test (CRISPRa for the three brakes, CRISPRi for the three drivers), peak condition, and a per-target go/no-go threshold, at an estimated ${\sim}96$ conditions over 3--4 months. Stating the falsification condition and the cost up front converts the nominations from a ranked list into a fundable, disprovable experiment.

\widefig{figs/fig_validation_bet.png}{0.98}{\textbf{A falsifiable prospective bet and a costed validation panel.} (a) The flagship SIK2 prediction with its explicit falsification condition. (b) The six-lead validation panel with directional test, peak condition, screen effect size, and powered sample size, plus a validation cost card. CRISPRa tests the three brakes (agonize); CRISPRi tests the three drivers (antagonize).\label{fig:valbet}}

\subsection{Generalisation to a second, patient-derived disease signature (multiple sclerosis)}
To test whether the framework is dataset-agnostic rather than tuned to our curated inflammatory program, we re-ran the entire directional pipeline against an independent, patient-derived multiple-sclerosis (MS) signature defined by disease-vs-healthy enrichment on an integrated CD4$^+$ atlas of 1{,}884{,}217 cells across two cohorts (an inflammation landscape, \emph{Nat Med} 2025; a circulating-CD4 landscape, Yasumizu, \emph{Cell Genomics} 2024), with MS present in both as a replication pair. The direction derives from a genome-wide disease-vs-healthy enrichment table (17{,}165 genes, per-cohort correlations $\rho_B/\rho_C$; released as \texttt{disease\_direction\_MS\_full.parquet}), from which our analysis takes the top-ranked subset. Two confounds had to be controlled first, both instances of the failure mode this paper critiques. \textbf{Signature contamination:} the raw disease direction was dominated by ambient/lineage genes (90\% of its top-10 were haemoglobin, platelet, immunoglobulin, myeloid or sex/mitochondrial transcripts) and driven by a single cohort (median $\rho_C/\rho_B=4.6$); removing these classes and requiring balanced cross-cohort replication recovered a direction of canonical MS CD4 biology (cytotoxic GNLY/PRF1, antigen presentation HLA-DPB1/NLRC5, Th1/Th17 CXCR3/TBX21/RORA, AP-1 BATF; Fig.~\ref{fig:msgen}a). \textbf{Pleiotropy:} raw reversal scores tracked perturbation footprint ($\rho=-0.33$ with the number of DE genes), so housekeeping knockdowns topped the naive ranking; a footprint-matched empirical null removed this dependence ($\rho=-0.07$). After both corrections, \textbf{no knockdown survives genome-wide FDR} --- the honest headline --- but the signal emerges where the thesis predicts, at the genetics intersection: the 21 MS genetic-anchored genes (Open Targets $+$ a CD4 cis-eQTL panel) are enriched for correctly-directed reversal (86\% vs 58\% background, one-sided binomial $p=0.007$; Mann--Whitney $p<10^{-5}$; Fig.~\ref{fig:msgen}b). The genetic-anchored druggable leads are IL2RB, TYK2, IL2RA, IL7R and TNFRSF1A (Fig.~\ref{fig:msgen}c), and the directional calls match approved-MS-drug mechanism in 5/5 measurable targets (natalizumab/ITGA4, fingolimod/S1PR1, alemtuzumab/CD52, TYK2 inhibitors, anti-IL2RA; Fig.~\ref{fig:msgen}d). This is a generalisation and validation result, not a new method: the same directional score, empirical null, genetics layer and druggability filter transfer unchanged to a second, patient-derived disease signature in the same cell type --- the framework applies to any CD4-T-cell-mediated disease with a definable signature, of which MS is one worked example.

\widefig{figs/fig_ms_generalization.png}{0.99}{\textbf{Generalisation to a patient-derived MS signature, confound-controlled.} (a) A footprint-matched empirical null removes the pleiotropy confound (reversal--footprint correlation $-0.33 \to -0.07$). (b) MS genetic-anchored genes are enriched for therapeutic-direction reversal (86\% vs 58\%; one-sided binomial $p=0.007$). (c) Genetic-anchored MS leads by footprint-matched effect size (teal = druggable; dashed = nominal $p<0.05$). (d) Corrected directional call matches approved-MS-drug mechanism in 5/5 measurable targets ($p=0.031$). No knockdown survives genome-wide FDR after correction; the signal is carried by the genetics$\times$direction intersection.\label{fig:msgen}}

\section{Discussion}
We have shown that a directional, context-aware reading of CD4$^+$ T-cell Perturb-seq produces target nominations whose \emph{direction} --- the therapeutically decisive property that genetics cannot supply --- externally replicates on matched cytokine readouts across two CRISPR modalities, and correctly fails on a mismatched readout. We deliberately report the negatives: our functional axis does not out-rank human genetics at recovering known targets, and no learned integration beats genetics alone. This is not a weakness of the framing but its scope: the contribution is an orthogonal, directional, calibrated layer, most valuable in combination with genetics rather than in competition with it. Three additions sharpen the translational reading. The directional call is internally self-consistent with how approved drugs act (\S\,\ref{fig:selfcons}), which is what licenses treating a knockdown screen as a direction predictor; a patient-anchored check shows the program elevated in tissue-accessible autoimmune diseases, though bulk-tissue context mismatch bounds this to an exploratory result; and \texttt{Path2Drug} turns an undruggable transcription-factor nomination into a druggable intermediate node on a mechanistic path, distinguishing pathway-grounded intervention points from study-biased network hubs. A head-to-head test against network propagation shows that the competing paradigm's apparent recovery advantage is largely a study-bias artefact --- its scores track node degree ($\rho=0.84$) --- whereas our functional readout is degree-independent by construction. Finally, the whole pipeline transfers to a second, patient-derived disease signature (multiple sclerosis): once signature contamination and perturbation-footprint pleiotropy are controlled, no knockdown survives genome-wide FDR, yet the MS genetic-anchored set is strongly enriched for correctly-directed reversal (86\% vs 58\%, $p=0.007$) and the direction matches approved-MS-drug mechanism in 5/5 measurable targets --- the paper's thesis in miniature on new data.

Several limitations bound the claims. The analysis re-uses a single published atlas; the strongest replication (Schmidt) shares laboratory lineage with the source data, so it is not fully arm's-length. The directional signal has modest resolution, and context-specificity is confounded by perturbation pleiotropy. The structural layer is a coarse co-folding filter, not a docking or free-energy calculation, and the nominations are computational hypotheses without wet-lab validation. The clearest next step is prospective testing --- arrayed knockdown with a cytokine readout --- of the highest-confidence directional nominations.

\subsection*{Data and code availability}
The 1{,}923-target directional shortlist (with calibrated direction probabilities), 11 predicted target--ligand complexes, per-nomination novelty-audit dossiers with citation-linked evidence, and all analysis code are released as project artifacts. The directional scorer, empirical-null test, calibration, and the directional gold-standard benchmark (including its honest negative result) are packaged as \texttt{directional\_screen}, a pip-installable module for reuse on any perturbation screen reducible to a perturbation\,$\times$\,readout matrix. The mechanistic path-tracing tool \texttt{Path2Drug}, its deterministic baseline comparison, and per-target outputs are released alongside. We emphasise this reusability as a first-class contribution: the method is dataset-agnostic and the CD4$^+$ T-cell atlas is its demonstration, not its limit --- the same directional scoring, empirical null, calibration, and path-tracing apply unchanged to any CRISPR or compound perturbation screen with a directional readout, in any cell type.

\begin{thebibliography}{9}
\small
\bibitem[Marson \textit{et al.}(2025)]{marson2025} Zhu~L., Dann~E. \textit{et al.} (2025) Genome-scale CRISPRi Perturb-seq in primary human CD4$^+$ T cells. \textit{bioRxiv}, doi:10.64898/2025.12.23.696273.
\bibitem[Schmidt \textit{et al.}(2022)]{schmidt2022} Schmidt~R. \textit{et al.} (2022) CRISPR activation and interference screens decode stimulation responses in primary human T cells. \textit{Science}, \textbf{375}, eabj4008.
\bibitem[Shifrut \textit{et al.}(2018)]{shifrut2018} Shifrut~E. \textit{et al.} (2018) Genome-wide CRISPR screens in primary human T cells reveal key regulators of immune function. \textit{Cell}, \textbf{175}, 1958--1971.
\bibitem[Passaro, Wohlwend \textit{et al.}(2025)]{boltz2} Passaro~S., Wohlwend~J. \textit{et al.} (2025) Boltz-2: accurate structure and binding-affinity prediction. \textit{github.com/jwohlwend/boltz}.
\end{thebibliography}

\end{document}
